% !TeX root = ../main.tex
\chapter*{preface}

数学系本科毕业之后, 决定在空闲时间整理一下以前的笔记. (纯闲得慌

其实本科刚开始的时候就热衷于 note taking 了 (也是闲得慌, 写日记心态). 不过为了让自己打字舒服一点一直在折腾, 导致很多笔记都很乱, 基本都是半成品. 
\begin{itemize}
    \item 一开始是手写. 但是只要你学数学几个月之后你基本就不会想再手写了. 因为证明基本全是文字, 打字比手写快太多, 而且 latex 熟了之后打一些希腊符号都比手写要舒服很多.
    \item 后面自然而然就学了 latex. 用 overleaf + elegantbook 模板记了很长一段时间. 但是由于手慢经常脑子里已经写了五句话手上只写了一句话, 搞得记笔记的时间有点太长了.
    \item 后面一年我基本是用本地环境 + vscode 的 latex 插件 (有一段时间用的 Cursor btw, 没啥区别), 配合 copilot 补全来记的. 这样舒服很多, 基本打字速度能赶上思考速度了. 但是 setting.json 一直没配完美, 实时编译速度还是太慢了. 
    \item 之后有一段时间我用的是 typst. typst 的实时编译是原生的, 比 latex 快很多. 记得这个时期是 gpt5 刚出来的时候, 那个以后 coding agent 已经很强大了. 我就让 ai 直接给我把一些常用末班都转成了 typst 模式并给我配好环境. typst 什么都好, 最大的问题是语言和 latex 不兼容, 而且社区小, 很难做写 paper 之类的正经的事情. 所以和 latex 混用的时候经常打混. 不过这段时间我发现折腾烦了, 其实记的比较少了. 
    \item opus 5.5 出来的时候, 某天我突发奇想能不能在 obsidian 里 vibe 一个 latex 编辑器. 然后这件事情居然被完成得非常完美. 此后我就不再折腾了. 
\end{itemize}
说这么多币话主要是在说明一件事, 就是之前我的 notes 在长期进展中是一个很乱的东西. 格式乱, 内容乱, 放的地方也比较乱, 平时也难抽空去管一下它们. 和我认识的朋友大概知道数学对我而言基本是纯兴趣爱好了. 第一年还有志于要做数学 research 但是后两年一直在干 swe (虽然这行业也是死得不能再死了, lol). 不过某天上班摸鱼的时候一拍脑袋决定应该翻新一下旧 notes 以及续写一些新 notes, 要不然之后可能就没兴趣继续学下去了. 然后我就开始在闲的时候抽空去整理它们. 现在正在写的这个 foundations of mathematics 的主题是我在大概把旧 notes 都整理成了统一格式的 latex 之后第一个新写的内容. \\

言归正传, 不说我自己的事情了. 这个 topics 的目的是捋一下数学作为一个形式逻辑体系本身的结构. 范围大概会包括基础的 first-order logic, proof theory, ZFC 体系公理集合论, 以及 type theory for lean4, 但是不涉及 category theory (我自己还没学明白,,). \\

引言到这里还没写完. 之后再补.

